\documentclass[11pt]{article}
\usepackage[margin=1in]{geometry}
\usepackage{amsmath,amssymb}
\usepackage[hidelinks]{hyperref}

\begin{document}

\section*{Human Review: Bounded Collapsed Almost Uncollapsed Map}

\noindent Original automatic proof: \texttt{solution\_packet.pdf}

\noindent Checked date: 15 June 2026

\noindent Status: Verified.

\noindent Verdict: The proof is correct. Under the printed wording of Problem 2.6, the circle map gives a valid solution.

\section*{Human Comments}

The source problem asks whether there is a map \(X\to Y\), for Banach spaces \(X\) and \(Y\), which is collapsed, almost uncollapsed, and bounded. The proposed example
\[
f:\mathbb R\to\mathbb R^2,\qquad f(x)=(\cos x,\sin x)
\]
does exactly this.

Boundedness is immediate, since the image lies in the Euclidean unit circle. The map is collapsed because it is periodic: for every \(t>0\), choose \(k\in\mathbb N\) with \(2\pi k\geq t\). Then
\[
|2\pi k-0|\geq t,\qquad f(2\pi k)=f(0),
\]
so the compression modulus satisfies
\[
\rho_f(t)=0
\]
for every \(t>0\). Hence \(f\) is not uncollapsed, i.e. it is collapsed in the terminology of the paper.

On the other hand, \(f\) is almost uncollapsed. If \(|x-y|=\pi\), then \(f(y)=-f(x)\), and therefore
\[
\|f(x)-f(y)\|_2=2.
\]
Thus the exact compression modulus satisfies
\[
\overline{\rho}_f(\pi)=2>0.
\]
So \(f\) is almost uncollapsed.

The construction also makes clear why the example works: periodicity produces arbitrarily far apart points with identical images, while the antipodal distance \(\pi\) gives uniform separation at one exact distance. In this sense, the example is not accidental; any suitably chosen periodic map with a uniformly separated exact distance would give the same phenomenon.

\section*{Infinite-Dimensional Variant}

The infinite-dimensional variant in the packet is also conceptually sound. One chooses a unit vector \(u\in X\), quotients by the cyclic subgroup \(2\pi\mathbb Z u\), and maps cosets to distinct unit vectors in an \(\ell_2\)-space. This recreates the same periodic-collapse mechanism along one direction while preserving separation at exact distance \(\pi\). This variant is not needed for the printed problem, since \(X=\mathbb R\) is already a Banach space, but it addresses any possible preference for infinite-dimensional domains.

\section*{Cocycle Comment}

Reading around Problem 2.6, I do not see any explicit requirement that the map be a cocycle or arise from an affine isometric action. The problem is stated for a ``map'' \(X\to Y\), and the finite-dimensional circle example is even continuous and Lipschitz. Therefore it answers the printed question.

In fact, the finite-dimensional example can also be written as a continuous bounded cocycle. Let \(R_t\) denote rotation by angle \(t\) on \(\mathbb R^2\), and define
\[
b(t)=R_t e_1-e_1=(\cos t-1,\sin t).
\]
Then \(b\) is a cocycle for the isometric representation \(t\mapsto R_t\), since
\[
b(s+t)=b(s)+R_s b(t).
\]
Subtracting the constant vector \(e_1\) does not change pairwise distances, so this cocycle has the same collapsed, almost uncollapsed, and bounded behaviour as the circle map.

The infinite-dimensional coset variant can also be written as an algebraic bounded cocycle. If \(A=2\pi\mathbb Z u\), \(I=X/A\), and \(Y=\ell_2(I)\), let \(X\) act on \(Y\) by permuting the basis vectors according to coset translation:
\[
\pi_x e_{[y]}=e_{[x+y]}.
\]
Then
\[
b(x)=\pi_x e_{[0]}-e_{[0]}=e_{[x]}-e_{[0]}
\]
is a bounded cocycle. It has the same collapse and exact-distance separation properties as the coset map. However, this infinite-dimensional cocycle is not continuous in general; arbitrarily small nonzero \(x\notin A\) already produce a jump in the image.

\section*{Review Outcome}

Accepted as a valid solution to Problem 2.6 as printed.

\end{document}
